\chapter{Introducción}
\label{sec:introduccion}

Actualmente, la seguridad en el desarrollo de software es fundamental. Sin
embargo, las malas prácticas y la falta de conocimiento siguen existiendo a
pesar del aumento de amenazas y la mejora continua de capacidades de los
atacantes.

Existe abundante literatura respecto a las vulnerabilidades con las que se puede
encontrar un desarrollador, como el OWASP Top 10~\cite{owasptop10}, que cataloga
las vulnerabilidades de software más críticas, incluyendo inyección, pérdida de
control de acceso, fallas criptográficas, entre otras. Sin embargo, la mejor
forma de conocer el riesgo real de las vulnerabilidades es entendiendo cómo
funcionan, para lo cual un enfoque práctico permite validar su correcto
entendimiento. Las investigaciones demuestran que el aprendizaje práctico es
significativamente más efectivo, con estudios mostrando que los participantes
retienen el 75\% de la información a través de práctica \textit{hands-on},
comparado con solo el 5\% mediante métodos teóricos
tradicionales~\cite{sanshandson}.

El problema a abordar en este trabajo es facilitar, tanto al equipo docente como
a los estudiantes, confirmar sus conocimientos teóricos a partir de una
plataforma automatizada que valide sus avances. Las y los desarrolladores
necesitan entornos controlados y seguros donde puedan experimentar y explorar
vulnerabilidades posibles dentro de una aplicación sin que esta sea de
producción. Para ello, se diseñó y desarrolló un sistema de evaluación
automático de ejercicios de ciberseguridad ofensiva.

Este proyecto se aproxima al problema desde tres perspectivas complementarias.
Desde lo \textbf{pedagógico}, se requiere una forma de evaluar el progreso que
capture la trayectoria completa del aprendizaje, no solo el resultado final. La
validación automática basada en el estado del sistema permite registrar cada
paso del estudiante durante el ataque, proporcionando trazabilidad real del
proceso de aprendizaje sin recurrir a soluciones artificiales como múltiples
\textit{flags} intermedias. Desde lo \textbf{técnico}, existe una oportunidad de
contribuir al ecosistema de herramientas \textit{open source} de detección de
intrusiones, publicando código, patrones y metodologías que permitan que la
comunidad académica evalúe y mejore estos enfoques. Desde lo \textbf{personal},
este proyecto representa la oportunidad de desarrollar experiencia práctica en
el diseño e implementación de sistemas complejos con múltiples componentes,
arquitecturas de capas bien definidas, y desafíos de observabilidad cercanos a
los que enfrenta la industria de ciberseguridad.

\section{Objetivos}

\subsection{Objetivo General}

Desarrollar una plataforma educativa autohospedable que permita crear entornos
de práctica de seguridad ofensiva personalizables, con validación automática del
progreso basada en estado de los sistemas comprometidos, facilitando el
aprendizaje práctico de vulnerabilidades sin depender de validaciones manuales.

\subsection{Objetivos Específicos}

\begin{enumerate}
    \item Diseñar e implementar una arquitectura modular que permita desplegar y
    gestionar múltiples entornos vulnerables simultáneamente mediante
    tecnologías de virtualización o contenerización, asegurando aislamiento y
    facilidad de configuración.
    
    \item Desarrollar un sistema de configuración declarativa que permita
    definir las condiciones de éxito de un ejercicio mediante archivos YAML,
    de modo que crear un ejercicio con validación automática requiera
    describir qué debe observarse y no programar la lógica de detección ni
    modificar la plataforma.

    \item Implementar el motor de validación automática que evalúe el estado de
    los sistemas comprometidos, determinando el progreso sin depender de
    \textit{flags} predefinidas, ya sea por archivos modificados, servicios
    comprometidos, usuarios creados, entre otros.

    \item Crear herramientas de seguimiento del progreso, acceso a enunciados de
    ejercicios y administración del sistema para educadores y estudiantes,
    implementadas inicialmente mediante línea de comandos.

    \item Desplegar la plataforma en la infraestructura del laboratorio de
    ciberseguridad CLCERT y verificar allí su operación completa: instalación,
    ejecución de extremo a extremo de los escenarios y acceso remoto de un
    participante.\footnote{Este objetivo, formulado originalmente como ``poner
    la plataforma a disposición del CLCERT para su evaluación en una instancia
    real de uso'', se precisa aquí con un criterio verificable.}
\end{enumerate}

\section{Evaluación}

La forma de evaluar este trabajo fue verificando la capacidad del sistema para
generar diferentes tipos de ejercicios con validación automática. El trabajo se
consideró exitoso si era posible crear ejercicios de seguridad ofensiva con
validación automática para al menos tres categorías entre: aplicaciones y
servicios web, bases de datos, \textit{pwning}, análisis forense e ingeniería
reversa. Se entendió por validación automática aquella que marca el ejercicio
como completado automáticamente cuando el participante logra explotar una
vulnerabilidad, acceder a un archivo específico, comprometer un servicio u
obtener privilegios de usuario, sin requerir la entrada manual de \textit{flags}
o códigos predefinidos.

\section{Solución Propuesta}

La solución consiste en una plataforma que permite a los educadores cargar
imágenes de sistemas vulnerables y definir reglas de validación mediante
archivos YAML, generando automáticamente entornos de práctica personalizados con
validación automática.

\subsection{Arquitectura General}

El sistema implementa una arquitectura dividida en cuatro capas con
responsabilidades claramente delimitadas. Esta arquitectura emergió del análisis
de pruebas de concepto~\cite{flexipwn-poc} desarrolladas durante el trabajo
previo al semestre, donde se identificó que la separación de responsabilidades
entre el sistema vulnerable, el mecanismo de detección y la lógica de evaluación
es fundamental para lograr flexibilidad y extensibilidad.

\begin{description}
    \item[Capa 1 — Sistemas aislados] Contiene todos los sistemas y redes que
    componen el escenario de ataque. La implementación actual los materializa
    como contenedores, aunque la arquitectura es agnóstica al origen de
    aislamiento y podría sostenerse a futuro sobre tecnologías de
    virtualización. Se divide en dos subcapas: la Capa 1a, el sistema vulnerable
    objetivo (el contenedor vulnerable), y la Capa 1b, el entorno atacante desde
    el cual el estudiante ejecuta el ataque (el contenedor atacante).

    \item[Capa 2 — Monitoreo pasivo] Observa pasivamente el estado de los
    contenedores con acceso de lectura externo al sistema. La única excepción a
    esa lectura externa pura es la captura de tráfico de red, que requiere un
    contenedor auxiliar (\textit{sidecar}) corriendo \texttt{tcpdump}, por ser
    la forma adecuada de capturar el tráfico del escenario. Los escenarios de
    referencia asumen tráfico en claro, sin TLS, para que la captura resulte
    legible. La responsabilidad de esta capa es solo generar un registro
    estructurado de eventos detectables: modificaciones del sistema de archivos,
    procesos creados, entradas de \textit{log} y tráfico o conexiones de red,
    entre
    otros. Su comportamiento es análogo al de un SIEM (Security Information and
    Event Management) acotado a un único escenario de ataque.

    \item[Capa 3 — Motor de evaluación] Consume los eventos generados por la
    Capa 2 y evalúa reglas declarativas definidas en archivos YAML que
    especifican las condiciones de victoria. Determina automáticamente cuándo un
    participante ha completado objetivos sin requerir la entrada manual de
    \textit{flags}.

    \item[Capa 4 — Administración] Gestiona el ciclo de vida completo de los
    entornos de práctica y mantiene registro del progreso de todos los
    participantes. Proporciona interfaces de línea de comandos principalmente
    para educadores, incluyendo vistas agregadas para seguir varios entornos
    en paralelo.
\end{description}

\subsection{Flujo principal}

\begin{enumerate}
    \item \textbf{Carga de escenarios}: el educador entrega una imagen del
    sistema vulnerable junto con un archivo YAML que define los objetivos del
    ejercicio, las condiciones de validación y los criterios de éxito.

    \item \textbf{Generación de entornos}: la plataforma instancia la imagen
    proporcionada, creando un entorno aislado por cada participante que
    intentará explotar la vulnerabilidad.

    \item \textbf{Validación automática}: el motor de validación examina
    constantemente el estado del contenedor, comparando el estado actual contra
    las reglas definidas en el archivo YAML.

    \item \textbf{Gestión}: tanto educadores como estudiantes utilizan
    herramientas de línea de comandos para interactuar con sus respectivos
    entornos y monitorear el avance.
\end{enumerate}

\subsection{Perfil de usuario esperado}
\label{sec:perfil-usuario}

FlexiPwn asume un educador con manejo de línea de comandos y familiaridad
básica con contenedores: debe poder instalar Docker en modo \textit{rootless}
una vez, construir o adaptar una imagen del sistema vulnerable y escribir un
archivo YAML. No es el tipo de usuario final que solo tiene que saber usar una aplicación web, sino
que corresponde al perfil habitual de quien dicta un curso de
seguridad ofensiva. Lo que la plataforma elimina es la necesidad de programar
el mecanismo de detección, es decir, configurar mecanismos para observar el
sistema objetivo, capturar y correlacionar eventos, y escribir el código que decide cuándo un objetivo se
cumple, que es la parte costosa de montar validación automática. Ese es el
sentido en que la plataforma resulta accesible: no porque rebaje los
requisitos de infraestructura, sino porque sustituye el trabajo de
instrumentación por una descripción declarativa.

Además, ese piso técnico podría considerarse el más bajo entre las alternativas
autohospedables analizadas en el Capítulo~\ref{sec:estado-del-arte}: CTFd
requiere desplegar un servicio web con base de datos, y los \textit{cyber
ranges} como Tectonic o CyberRangeCZ se operan con herramientas de
infraestructura como código, entre ellas Packer, Terraform y Ansible, sobre
proveedores de nube o hipervisores.

\section{Declaración de uso de inteligencia artificial}

En la elaboración de esta memoria se utilizaron herramientas de inteligencia
artificial generativa, concretamente los modelos Claude de Anthropic (familias
Opus y Sonnet 4.x), empleados principalmente a través de la aplicación de
escritorio en su modo Cowork. Estas herramientas se usaron a lo largo de todo
el trabajo: en la revisión de literatura y la comparación de
plataformas, en la discusión de decisiones de arquitectura, 
en la generación y refactorización de código, en la redacción y
edición de este documento,  y en tareas de apoyo
como el control de versiones y la organización del proyecto. Todo el contenido
generado o asistido fue revisado, evaluado y validado por el autor,
quien asume plena responsabilidad sobre el trabajo presentado. La
inteligencia artificial se empleó como herramienta de apoyo y no se reconoce
como coautoría.

El modo Cowork se utilizó como un acelerador del trabajo: permitió delegar la
realización del trabajo, como escribir código, redactar y editar, ejecutar y
verificar, conservando el rol del autor como arquitecto de las decisiones.
Trabajar de esta forma permitió mantener la constancia durante todo el semestre,
compatibilizando con las labores del autor como desarrollador de aplicaciones
móviles. La integración de la herramienta estuvo acompañado de la verificación:
la relectura crítica de cada resultado, la comparación de las afirmaciones y
planeamientos contra el código fuente, y el uso de una memoria persistente entre
sesiones para mantener el contexto y evitar contradicciones. El criterio
principal fue no dar por válido ningún resultado por el solo hecho de que
``funciona'' a primera vista.

Los prompts siguieron unos patrones recurrentes. Para una tarea de inicio a fin,
se partía entregando el contexto (usualmente un resumen de lo trabajado en
sesiones anteriores), seguido de lo que había que hacer y de los puntos que
valía la pena discutir, pidiendo siempre alternativas para elegir la
implementación más adecuada. Se cerraba con algún criterio de validación y con
notas que precisaban la ejecución. A mitad de una tarea el patrón cambiaba: se
revisaba lo ejecutado y se respondía con un punteo extenso de dudas, cambios
deseados y aspectos por considerar, y solo se confirmaba el avance al paso
siguiente cuando el resultado era satisfactorio. La mayoría de veces, se
recordaba al modelo lo que debería tener en memoria, para asegurar un correcto
alineamiento de la tarea a realizar. Los prompts breves se usaban para algunos
casos puntuales cuya discusión convenía expandir, y al final de una sesión
cargada se solicitaban los cierres habituales: formateo, guardado en memoria,
mensaje de \textit{commit} si se requería y un prompt de traspaso que orientara
a la sesión siguiente.

A modo ilustrativo, dos ejemplos representativos del tipo de instrucción
entregada:

\begin{itemize}
    \item Un prompt que funcionó bien: ``Retoma el informe: primero una foto del
    estado del repositorio con los commits nuevos. Luego valida afirmación por
    afirmación contra el código, marcando cada una como verdadera, falsa o
    imprecisa con su archivo y línea. Entrega un reporte ordenado por criticidad
    antes de editar y espera mi confirmación. Reglas: 80 caracteres, prosa en
    presente y opciones ante cada alternativa.'' Funcionó porque aportaba
    contexto, descomponía la tarea, exigía validación y opciones, y fijaba
    restricciones explícitas.
    \item Un prompt que funcionó mal: una instrucción amplia y sin contexto como
    ``mejora el capítulo de implementación'', que al no acotar el objetivo ni
    los criterios de aceptación producía cambios genéricos o de alcance excesivo
    que luego exigían una corrección considerable. La lección es que los prompts
    pobres en contexto y sin criterios claros rendían resultados de bajo valor.
\end{itemize}

Surgió un momento interesante en la etapa final del trabajo, al estudiar una
prueba de concepto de escape de contenedores para comprender el funcionamiento
de este tipo de \textit{exploit} y respaldar el análisis de contención del
entorno. Al
tratar un tema sensible, las primeras consultas sobre una prueba de concepto
realista del CVE~\cite{cve-runc-2025, poc-cve-2025-31133} en cuestión fueron
rechazadas por las políticas de uso de la herramienta. Tras solicitar al equipo
de Anthropic un acceso ampliado y explicar su propósito educativo dentro de este
trabajo de ciberseguridad, la consulta pudo retomarse. A partir de ahí, el
modelo resultó de gran ayuda para seguir el \textit{exploit} paso a paso,
consiguiendo
compararlo con otros CVE relacionados, distinguiendo, por ejemplo, los que se
ejecutan desde el interior del contenedor de los que afectan directamente al
\textit{host}.

Si el trabajo se reiniciara, se cambiarían algunas formas del uso de la
inteligencia artificial. Se establecería antes una rutina de revisión y
verificación sistemática, para no depender de que un resultado ``parezca''
correcto. Además, se sería más disciplinado con la memoria persistente y los
traspasos entre sesiones, que resultaron clave para mantener la coherencia a lo
largo de un trabajo asistido por IA.
